\section{Related Work}
\label{sec:related}

Which ingredient is new? None. The handshake is.

\textbf{Language models for operations.} InventoryBench established the
complementarity this paper starts from, and its four strategies are our arms 1,
3, 5 and 11~\cite{inventorybench}. A design-time session can synthesize reusable
operations research algorithms outright~\cite{llm-or-algorithms}, and InvEvolve
evolves white-box inventory policies online~\cite{invevolve}. Our model does the
opposite of both: it estimates the environment rather than the policy, writes no
code, and every action stays fixed-form operations research. InvEvolve has no
public code artifact as of our search, so any comparison against it must be
labeled a faithful reimplementation, and we have not implemented one. Language
models reach the operations loop in other shapes as well: as a multi-agent
inventory manager~\cite{invagent}, as a reader of event text for demand
forecasting~\cite{eventcast}, and as a human-in-the-loop instruction channel
into a controller~\cite{instructmpc}; the same replenishment setting has a
multi-agent reinforcement-learning benchmark~\cite{mabim}, and text-game agent
evaluation has its own benchmark~\cite{textarena}. Field evidence supports the
interface shape: judgmental adjustment of a forecast input
raised profitability by 4.92 percent on average in a retail
experiment~\cite{kesavan-field}, and average United States import shipping delay
rose by 21 days between 2018 and 2024~\cite{carreras-valle-delays}.

\textbf{When to call, and what to commit to.} Event-triggered invocation with
threshold, CUSUM, sequential probability ratio~\cite{wald-sprt},
Bayesian~\cite{bocpd} and learned rules is already
established~\cite{event-triggered-llm}. Agents ask the same when-to-call
question of themselves, answered by gating over fast-slow planning~\cite{asscg},
budgeted replanning~\cite{brace} and learned model routing~\cite{routellm}, so
we treat invocation as fixed infrastructure and claim nothing about it. Structuring a plan as falsifiable commitments with confirming
and falsifying evidence, and running a hypothesis lifecycle with provenance and
belief updates, also exist already~\cite{fcpagent, hep}. What differs here is
that the commitment parameterizes an \emph{exogenous stochastic model} whose
influence is authorized by a preregistered sequential statistic and compiled into
inventory-control inputs, so the object is tested rather than merely recorded.

\textbf{Statistics and inventory theory.} We use e-values and anytime-valid
inference as given~\cite{vovk-wang-evalues, ramdas-savi, grunwald-safe-testing},
with the classical sequential and changepoint machinery behind them~\cite{ville,
wald-sprt, page-cusum, lorden, hinkley-changepoint}. The
nearest construction is the e-detector, which sums
processes started at every time and controls the average run
length~\cite{e-detectors}, while we start one only at model-selected stopping
times and control per-episode false activation by budget splitting. A theorem-backed
\emph{retirement} rule would need confidence sequences~\cite{howard-confseq},
specifically the forward and backward confidence-sequence intersection
test~\cite{shekhar-backward-cs}, which we do not adopt, which is why
our retirement rules are labeled empirical. The controller is the textbook
one~\cite{clark-scarf, zipkin-foundations}, with base-stock heuristics for
lost-sales systems, censored demand included, long
established~\cite{zipkin-lost-sales, huh-lost-sales}, and our transit-pause
semantics follow an established simulation convention in which material already
moving stops progressing~\cite{stockpyl}.
